% June Week 5 weekly: Mon 2026-06-29 to Sun 2026-07-05.
% This week: leader-facing research summary of v1.8.0 outcome, public
% long-context baseline fact-base, and v2.0.0 decision options.

\section{June Week 5 (Mon 2026-06-29 to Sun 2026-07-05)}

\begin{frame}{Reporting Period}
  \begin{center}
    \textbf{Plan 08 / Paper B --- June Week 5} \\[0.3em]
    \small Mon 2026-06-29 to Sun 2026-07-05 \\
    \scriptsize Leader update: research outcome, long-context facts, v2.0.0 options, and resource ask.
  \end{center}
\end{frame}

\begin{frame}{Executive Message}
  \begin{block}{From ``make compression work'' to a long-context research thesis}
    The v1.8.0 work showed that a learned memory plus a gate can be useful, but
    it also showed that the carrier itself is not the whole story. This week we
    shifted to a public long-context fact-base: when do existing black-box
    methods work, when do they fail, and what gap should the next paper target?
  \end{block}
  \begin{itemize}
    \item \textbf{Outcome so far:} v1.8.0 has a working memory+gate recipe on high-headroom tasks.
    \item \textbf{New research lesson:} no fixed long-context baseline wins across length, task, and model architecture.
    \item \textbf{Decision needed:} whether to invest in a v2.0.0 paper around adaptive, self-verified evidence memory.
  \end{itemize}
\end{frame}

\begin{frame}{What We Did This Week: Public Long-Context Fact-Base}
  \small
  \begin{itemize}
    \item Ran / organized a broad black-box baseline campaign on public long-context benchmarks.
    \item Covered window truncation, KV eviction, prompt compression, RAG, full-context, and no-context references.
    \item Built a failure-mode map instead of only reporting a leaderboard.
    \item Wrote the fact-base and diagnosis notes that define the next paper's gap.
  \end{itemize}
  \begin{block}{Core output}
    The research artifact is an empirical map: method $\times$ length $\times$
    task type $\times$ architecture. It tells us where fixed long-context
    coping mechanisms break.
  \end{block}
\end{frame}

\begin{frame}{Problem Statement: Long Context Is Not One Problem}
  \small
  \begin{itemize}
    \item Sometimes full context is needed; sometimes it is expensive; sometimes it actively hurts.
    \item Sometimes a recent window is enough; sometimes the answer is far away.
    \item Sometimes retrieval works; sometimes there is no lexical overlap.
    \item Sometimes KV eviction is available; sometimes the model architecture has no KV lever.
  \end{itemize}
  \begin{block}{Goal}
    Stop treating long context as a single compression problem. Identify the
    regimes, then design a method that adapts to the regime.
  \end{block}
\end{frame}

\begin{frame}{v1.8.0 Outcome: What Worked}
  \small
  \begin{itemize}
    \item Fixed earlier training/eval artifacts: termination supervision and joint answer signal recovered real compression behavior.
    \item On high-headroom tool-use, the memory path can be strong: bfcl-style results recovered to the expected range.
    \item The gate is valuable because it can route away from unsafe memory and preserve a safer path.
    \item Flip-one ablations show the carrier is not fragile: many knobs move compress by only a few points.
  \end{itemize}
  \begin{block}{Interpretation}
    v1.8.0 validates the direction, but it also tells us that polishing the
    compressor is not enough. The contribution has to be the decision layer:
    when to trust memory and when to fall back.
  \end{block}
\end{frame}

\begin{frame}{v1.8.0 Caveats: What We Should Not Overclaim}
  \small
  \begin{itemize}
    \item ``Gated $\ge$ full'' can be tautological if the fallback is defined as full; the real claim needs held-out or conformal calibration.
    \item Best gate signal is benchmark-dependent in the current recipe; this is inelegant.
    \item Some tasks have little headroom, so compression cannot create value there.
    \item Cross-domain and long-window behavior still need a cleaner public benchmark story.
  \end{itemize}
  \begin{block}{Implication}
    v1.8.0 should be treated as the working prototype. v2.0.0 should sharpen the
    claim into a self-verified, out-of-sample reliability story.
  \end{block}
\end{frame}

\begin{frame}{Fact-Base Question}
  \small
  \begin{block}{Question}
    Which long-context coping method works, when, and why?
  \end{block}
  \begin{itemize}
    \item Baselines: window / truncation, KV eviction, prompt compression, RAG,
      full-context and no-context references.
    \item The goal is not to crown one winner, but to map failure modes across
      length, task type, and model architecture.
    \item All paper-facing claims should be grounded in public benchmarks.
  \end{itemize}
\end{frame}

\begin{frame}{Key Long-Context Facts}
  \small
  \begin{itemize}
    \item \textbf{Attention-based KV eviction can collapse past 8k--16k} because the observation window mis-ranks far evidence.
    \item \textbf{Attention-free or reconstruction-style eviction is more length-robust on retrieval,} but not universally semantic-robust.
    \item \textbf{Ranking flips by task type:} retrieval, extractive QA, global reasoning, and literary MC behave differently.
    \item \textbf{Full context can hurt} on some tasks with distractor-heavy or capability-limited settings.
    \item \textbf{RAG is strong with lexical overlap,} but weak for abstractive or global tasks.
    \item \textbf{Linear-attention / GDN-style models do not expose a KV lever.}
  \end{itemize}
  \begin{block}{Conclusion}
    A fixed context-reduction method is not enough.
  \end{block}
\end{frame}

\begin{frame}{v2.0.0 Paper Direction}
  \small
  \begin{block}{High-level thesis}
    Long-context reasoning needs adaptive evidence memory,
    not one fixed compression or retrieval rule.
  \end{block}
  \begin{itemize}
    \item Full context can be too expensive or misleading.
    \item Truncation only works when the answer is local.
    \item RAG only works when lexical overlap is available.
    \item KV eviction is model-architecture dependent and can fail at long length.
  \end{itemize}
  \begin{block}{Method framing}
    Build a memory that preserves evidence, verify whether it preserved enough,
    and route to the cheapest reliable path.
  \end{block}
\end{frame}

\begin{frame}{v2.0.0 Ideas}
  \small
  \begin{enumerate}
    \item \textbf{Self-verified evidence memory:} train the memory to preserve evidence;
      use that preservation/reconstruction signal as the reliability check.
    \item \textbf{Reliability-gated fallback:} use extracted knowledge when safe;
      otherwise fall back to raw context, retrieval, or a higher-cost path.
    \item \textbf{Hybrid evidence memory:} combine a raw recent window with compressed global memory.
    \item \textbf{Adaptive budget:} easy contexts get fewer memory tokens; dense or ambiguous contexts get more.
  \end{enumerate}
\end{frame}

\begin{frame}{What We Need To Run Next}
  \small
  \begin{itemize}
    \item \textbf{M1: Hybrid Evidence Memory} --- reader sees raw recent window + compressed global memory.
    \item \textbf{M2: Conformal router} --- held-out reliability calibration instead of in-sample threshold picking.
    \item \textbf{M3: Long-context necessity wave} --- semantic-spread tasks plus retrieval honest-negative.
    \item \textbf{M4: Transfer matrix} --- train on one domain, test on another; measure whether memory is domain-invariant.
    \item \textbf{M5: Figures} --- turn the fact-base into length, task, and method-failure figures.
  \end{itemize}
\end{frame}

\begin{frame}{What Leaders Need to Decide}
  \small
  \begin{description}
    \item[Option A: Fact-base paper] Publish the long-context baseline map first; method comes later.
    \item[Option B: Method paper] Build v2.0.0 now around self-verified evidence memory + conformal routing.
    \item[Option C: Two-stage plan] Fact-base as Section 1 / motivation, then v2.0.0 as the method contribution.
  \end{description}
  \begin{block}{Recommendation}
    Option C: use the fact-base to define the gap, then spend compute only on the
    minimum v2.0.0 experiments needed to show the gap can be closed.
  \end{block}
\end{frame}

\begin{frame}{Recommended Next Step}
  \small
  \begin{itemize}
    \item Freeze the public long-context fact-base and produce the figures.
    \item Implement Hybrid Evidence Memory: raw window + compressed memory.
    \item Add conformal reliability routing and report risk--coverage / cost--coverage.
    \item Run the 5x5 transfer matrix to test cross-domain memory generalization.
    \item Decide whether v2.0.0 is a full method paper or a fact-base + position paper.
  \end{itemize}
  \begin{block}{Ask}
    Approve GPU budget for the v2.0.0 public benchmark experiments, especially
    the hybrid-memory and transfer-matrix runs.
  \end{block}
\end{frame}
